% Training-set size: the same ablation for every task, so a model paper
% inputs this file inside its ablation section and then reports its own
% findings. It appears once the task has the runs.
% The prose follows docs/writing.md.
\ifgames
\subsection{Training-set size}
\label{sec:training-size}

To estimate how much a model's accuracy depends on the amount of training
data, we retrain it on a half and on a quarter of the training games and
evaluate it on the full test set (Table~\ref{tab:games},
Figure~\ref{fig:games}). The subsets are selected by a hash of the game
identifier and are nested, so the quarter is contained in the half. A network
is given twice the epochs on half the games and four times the epochs on a
quarter, so the number of optimisation steps is constant and only the amount
of data varies.
\ifmore
We also train on more data. A larger sample of the same source is split in
the same way as the standard dataset, so the validation and test sets are
unchanged and no test game appears in the larger training set. A network
trained on a larger set is given proportionally fewer epochs, so that it
again takes the same number of optimisation steps as its base. Unlike a model
trained on a subset, a model trained on more data is a candidate model, and
the acceptance criterion applies to it.
\fi

\begin{table}[htbp]\centering\footnotesize
\block{study-games}
\caption{Training-set size. Each model is trained on a fraction of the
training games\ifmore, or on a multiple of them from the larger sample,\fi{}
and evaluated on the full test set. The gain is paired with the same model
trained on the standard training set. A negative gain means that the model is
worse, and bold marks statistical significance.}
\label{tab:games}
\end{table}

\begin{figure}[htbp]\centering
\block{study-games-figure}
\caption{Test score as a function of the number of training games. Each step
to the right doubles the training set.}
\label{fig:games}
\end{figure}
\fi
